% Generated from changelog.md by changelog_to_tex.py -- do not edit by hand.
\documentclass[11pt]{article}
\usepackage[margin=1in]{geometry}
\usepackage[T1]{fontenc}
\usepackage{lmodern}
\usepackage{parskip}
\usepackage{enumitem}
\usepackage[hidelinks]{hyperref}
\setlist{nosep,leftmargin=1.4em}
\begin{document}
\title{PyATPG development log}
\author{}
\date{}
\maketitle
What changed, why, and what it measurably did. Newest first. The LaTeX version (\texttt{docs/changelog.tex}) is generated from this file by \texttt{python docs/changelog\_to\_tex.py}; edit this file, not the \texttt{.tex}.

Numbers marked \emph{detector only} come from running the trigger detector on the full-scan model with no ATPG and no LLM (a few seconds per circuit). Numbers from full runs say so.

\section*{2026-10-11: combinational triggers, LLM selection on a heuristic floor}
Prompted by the first held-out runs (compare mode, full scan, LLM on): s35932-T100/T200/T300, s38417-T100/T200/T300, s38584-T100/T200/T300. s38584 and s35932-T100 were found whole; the rest showed four problems. Recall with / without the LLM:

\begin{itemize}
  \item s35932-T200: 0.20 / 0.20, verdict inconclusive.
  \item s38417-T100: 0.91 / 0.91, Trojan suspected (0.78).
  \item s38417-T200: 0.36 / 0.71, Trojan suspected (0.78).
\end{itemize}

\subsection*{Combinational trigger trees (\texttt{trigger\_detector.py})}
\begin{itemize}
  \item \textbf{Problem.} s35932-T200 and s38417-T100/T200 have no trigger register: Tj\_OUT1..8 feed only Tj\_OUT1234/5678, which feed only Tj\_Trigger, which drives the payload gates. Every strong rule needs a register or an output splice. On s35932-T200 the leaves are 1 in \textasciitilde{}40\% of random full-scan patterns, so no rarity sign sees the tree; on s38417-T200 Tj\_OUT5 toggles 4 times in 16384 patterns and split the never-toggling cluster. The payload gates (\texttt{n4263\_Tj\_Payload}, \texttt{n4122/4149/4160/4161\_Tj\_Payload}, probability 0.06 to 0.44) were in no list at all.
  \item \textbf{Check before changing it.} A fanout-free tree of rare gates is not evidence: clean ISCAS s38584 has a 19-gate one with a root that never toggles, s838 several 11-gate ones. What the three trees share is the hand edit inside them, \texttt{NOR(Tj\_OUT1, Tj\_OUT1, ...)}. Duplicate-input gates do occur in clean ISCAS-85 (5 gates over c1908, c2670, c3540), so the edit alone stays a lead.
  \item \textbf{Change.} New evidence \texttt{trigger\_tree} (0.70) and \texttt{trigger\_payload} (0.70, payload family). From a duplicate-input gate, follow the single consumer while it is a conjunction over it (its rare value needs this gate's rare value; an XOR or OR payload is not), to a root rare in under 5\% of patterns. The root and the fanout-free tree under it (at least 5 gates, at most 3 levels) are the trigger; the gates the root drives are the payload, and their other inputs are host signals the region never grows into.
  \item \textbf{Effect (detector only, synthetic).} The Trust-Hub netlists for s35932/s38417 are not in the repository, so the rule was checked on their shape inserted into clean ISCAS s38417, s35932, s15850 and s38584 (rare and common leaves, 1 or 4 payload gates): every Trojan gate flagged strong, no host gate flagged strong. Without the duplicate input nothing is found, which is the rule's limit. Unchanged: clean s15850, s38417, s35932, s38584 (ISCAS and Trust-Hub TjFree), the clean ISCAS-85 circuits with duplicate inputs (nothing strong), s38584-T300 TjIn/TjIn\_fixed (identical output), and all 120 generated Trojan netlists under \texttt{atpg testing/c*}.
\end{itemize}

\subsection*{LLM selection: heuristic floor, same budget, every gate scored (\texttt{targeting.py})}
\begin{itemize}
  \item \textbf{Problem.} On s38417-T200 recall \emph{with} the LLM was 0.36 against 0.71 without, for three reasons visible in the report.
  \item \textbf{Cause 1: the model's picks replaced the heuristic's.} The heuristic ranks all 9386 gates; the model saw 306 (3.3\%), and Tj\_OUT5..8 were not among what it kept.
  \item \textbf{Cause 2: the two lists were not the same size.} 163 gates with the LLM, 204 without (T100: 131 vs 206), so part of the gap was budget, not choice.
  \item \textbf{Cause 3: each batch kept only its best 34 gates for the next.} A gate dropped there was gone, and the report did not record the pool, so which of the two happened could not be told.
  \item \textbf{Missed signal.} The circuit-analysis model, which sees the 120 rarest gates, named Tj\_OUT5678, Tj\_OUT7 and Tj\_OUT8 as a suspicious region in the same run, but it ran after ATPG and fed nothing.
  \item \textbf{Change: same budget, heuristic floor.} Both modes test the same number of gates (the existing 30\% / 20-200 budget, or the new \texttt{target\_gates} setting). LLM runs keep the heuristic's top picks and the model fills the remaining \texttt{llm\_share} of the list (new setting, default 0.5); slots it leaves empty (no gate scored 30+) go to the heuristic's next picks. An LLM run can lose at most the heuristic's last picks, and \texttt{llm\_only\_hits} is exactly what the model added.
  \item \textbf{Change: every gate scored.} The model scores every gate it is shown (0-100), in independent batches with no trimming between them; scores are ranked globally, and gates left out of an answer are asked for once more. It is shown only gates outside the floor and the forced detector findings.
  \item \textbf{Change: circuit analysis first.} It now runs before selection, and the gates of its suspicious regions of at most 25 gates take the model's slots first (it often returns one region spanning its whole 120-gate input, which says nothing). The no-LLM list never gets them.
  \item \textbf{Change: logged.} The report's \texttt{targeting.context} holds the pool, every score, the floor, the model's picks and the heuristic fill; the ATPG page shows how many gates the model scored and how its slots were filled.
  \item \textbf{Effect (fake model, s38584-T300 TjIn).} 444 of 13,010 gates scored in 4 requests of 111 gates under the 7000-token budget (306 seen before); both lists 200 gates with the same 100-gate floor. Not yet measured with a real model.
\end{itemize}

\subsection*{Ground truth leaves out undriven nets (\texttt{trojan\_metrics.py})}
\begin{itemize}
  \item s38584-T300 ships its trigger output \texttt{Trigger\_out} undriven; it was promoted to a free input and counted as a Trojan gate no rule could find. T300 TjIn ground truth: 1937 $\rightarrow$ 1936 gates.
  \item s35932-T300 had no ground truth at all: it reads undriven \texttt{TjTrigger} and \texttt{Stage3}, and no net in the parsed netlist has a Trojan name. The Trojan's own module was not in the upload (or drives no output and was removed as dead logic); the run needs the whole Trust-Hub folder.
\end{itemize}

\subsection*{Verdict suspects (\texttt{verdict.py})}
\begin{itemize}
  \item XOR-shape-only matches are left out of the suspect list when the trigger detector has a strong finding: on the s35932/s38417/s38584 runs they hit no Trojan gate in 7 of 7 scored runs, and up to 232 of them diluted the list (s38417-T100 verdict-suspect precision 0.04).
  \item The "inconclusive" summary said every lead was an XOR-splice match even when the leads were never-toggling groups or duplicate inputs; it now counts each kind.
\end{itemize}

\subsection*{Not fixed}
\begin{itemize}
  \item s38417 keeps 16 to 36 targeted faults aborted on host nets (\texttt{n2107}, \texttt{n2005}, ...) even at \texttt{-b 10000 -L}, holding targeted coverage at 85-89\%. None are Trojan faults.
  \item The s38584-T200 rerun selected 385 gates against 389 on 9 Oct (13 and 9 gates differ), precision 0.67 both times, so the "\textasciitilde{}115 fewer non-Trojan gates" expected from the XOR-shape change on 9 Oct did not appear. Whether the 9 Oct run was already on that code is not recorded in the report (\texttt{runs.backend\_commit} in the run database has it).
\end{itemize}

\subsection*{Next runs (none of today's changes is measured yet)}
\begin{itemize}
  \item Rerun s35932-T200 and s38417-T100/T200 (compare mode, full scan, LLM on). These now count as tuned circuits: report their numbers before and after this entry, not as held-out results.
  \item Re-upload s35932-T300 as the whole Trust-Hub folder.
  \item Run the Trojan-free versions (s35932, s38417, s38584 TjFree) to confirm the new rule raises no false alarm on the real netlists.
  \item Recall at several budgets: the same circuits at \texttt{target\_gates} 50, 100 and 200.
  \item Repeat runs for run-to-run spread need \texttt{LLM\_RESULT\_CACHE=0} on the backend; otherwise every repeat reuses the first run's cached scores.
  \item Round 2 (\texttt{detector\_in\_targeting} off) after the runs above.
\end{itemize}

\section*{2026-10-10: smaller report downloads}
Prompted by the s38584-T300 report download: 14.8 MB, too large to analyse.

\subsection*{Suspects reference vectors by index (\texttt{verdict.py})}
\begin{itemize}
  \item \textbf{Problem.} Every one of the 1998 verdict suspects carried the same 3 test patterns as full 1601-bit strings: 10.5 MB of the 14.8 MB file.
  \item \textbf{Change.} \texttt{activating\_vectors} now holds indices into the run's \texttt{test\_vectors.vectors}, where each pattern is stored once. Older reports with strings still load (the run DB only counts them).
  \item \textbf{Effect.} The same report, re-encoded: 14.8 MB $\rightarrow$ 5.2 MB.
\end{itemize}

\subsection*{Summary export (\texttt{Reports.tsx})}
\begin{itemize}
  \item \textbf{Problem.} Even at 5 MB the full report is too big to read or paste into an LLM, and most of it is long gate lists repeated across sections.
  \item \textbf{Change.} New \emph{Download Summary} button next to the full export. Lists longer than 25 keep their first 25 entries plus a total count (suspects are sorted by confidence, so the head is the strongest); test-pattern bit strings become \texttt{<N-bit pattern>}. The full export is unchanged, for the aggregation scripts.
  \item \textbf{Effect.} s38584-T300: \textasciitilde{}140 KB summary.
\end{itemize}

\section*{2026-10-09: payload muxes, measuring the LLM alone, fewer junk targets}
Prompted by the s38584-T200 report (compare mode, full scan, LLM on): recall 258/260, verdict \emph{Trojan suspected} (0.97), no LLM-only hits, targeted precision 0.48, and 347 of the 419 redundant faults inside the Trojan.

\subsection*{Payload multiplexer rule (\texttt{trigger\_detector.py})}
\begin{itemize}
  \item \textbf{Problem.} T200 and T300 both replace output \texttt{g10122} with a 2:1 mux whose select is the trigger. The parser expands a mux to NOT(s), AND(d0, NOT s), AND(d1, s), OR. The 3-gate payload walk from the trigger registers reached only part of it, so 2 (T200) to 3 (T300 fixed) mux gates were always missed.
  \item \textbf{Change.} New evidence \texttt{payload\_mux} (0.80, payload family): a mux of that shape whose select is already a strong suspect is flagged whole. A register's own load mux (select = its clock net) is excluded; without that exclusion the rule matched 1024 load-mux gates on T300.
  \item \textbf{Effect (detector only).} T300 TjIn\_fixed 2166 $\rightarrow$ 2169 of 2170 (the last miss is \texttt{g10122\_Tj\_\_ppo}, the host's original output signal). s15850-T100 precision 0.78 $\rightarrow$ 0.80. Shipped T300 TjIn unchanged: its select is the undriven \texttt{Trigger\_out}, a free input, so nothing upstream is flagged.
\end{itemize}

\subsection*{\texttt{detector\_in\_targeting} setting: measure the LLM on its own}
\begin{itemize}
  \item \textbf{Problem.} With the detector forcing \textasciitilde{}260 (T200) to \textasciitilde{}2000 (T300) gates into the fault list, the LLM could add at most the 2 to 5 gates left over. \texttt{llm\_only\_hits} was always empty and said nothing about the LLM.
  \item \textbf{Change.} Live setting \texttt{detector\_in\_targeting} (env \texttt{PYATPG\_DETECTOR\_IN\_TARGETING}, default on). Off: no detector gate is forced into the pool or the list, no marks appear in the LLM's table, the no-LLM baseline is the plain heuristic, and the run notes say so. The verdict still uses the detector. The LLM cache key includes the number of forced gates, so an off run never reuses an on run's answer.
\end{itemize}

\subsection*{XOR-shape matches forced in only when nothing strong was found}
\begin{itemize}
  \item \textbf{Problem.} The older XOR-splice detector forced its matches into every targeted list. On T200 19 of its 134 gates were Trojan gates; on T300 181 of 295. Every one of those was also a trigger-detector finding.
  \item \textbf{Check before changing it.} On the tool's own generated Trojans (c432, c1355, c2670, 18 circuits) the XOR shape finds 16 of 18 Trojan gates and the trigger detector finds 1. Dropping it everywhere would have wrecked those benchmarks.
  \item \textbf{Change.} Its matches are still reported and still verdict leads, but join the fault list only when the trigger detector has no finding of strength 0.7 or more. Expected effect on T200: the targeted list loses \textasciitilde{}115 non-Trojan gates and no Trojan gates (not yet confirmed by a full run).
\end{itemize}

\subsection*{T200: why its trigger can never fire}
\begin{itemize}
  \item ATPG proved \texttt{TrojanTrigger\_\_n37 /0} redundant in both runs: the net clocking the Trojan counter can never be 1, so the counter never loads.
  \item It is \textbf{not} the scan enable held at 0: none of the trigger's input nets (\texttt{n2601}, \texttt{n5111}, \texttt{n1422}, \texttt{n3279}, \texttt{n3121}, \texttt{n3180}) read \texttt{test\_se} (checked on the shared s38584 host netlist).
  \item Likely cause: \texttt{n3121} and \texttt{n3180}, two NANDs sharing input \texttt{n1288}, were never 0 together in 65,536 random full-scan patterns, while every other combination appeared. Not proven exhaustively (35-input cone); Atalanta's redundancy result is the proof that the clock is unreachable.
\end{itemize}

\section*{2026-10-09: dead triggers found structurally (\texttt{9ff5dc4})}
\begin{itemize}
  \item The "trigger never fires" flag matched net names containing "trig", so T200's clock net \texttt{TrojanTrigger\_\_n37} was missed and the report said \texttt{trigger\_never\_fires: false}.
  \item Now any stem stuck-at-0 on a net that clocks a logic-clocked register counts (the parser models it as Q' = clk ? D : Q, so clk = 1 is observable). The run notes report it with or without ground truth. Replayed on the T200 report: \texttt{true} for both runs.
\end{itemize}

\section*{2026-10-09: circuit-analysis prompt anonymised (\texttt{23f4f81})}
\begin{itemize}
  \item The circuit-analysis LLM stage sent real gate ids, the file name and 3000 characters of Verilog. Trust-Hub names its Trojan logic; on T300 the model reported a "Trojan payload carry chain".
  \item It now gets \texttt{g0..gN} aliases and statistics only; aliases in its answer are mapped back to real ids. On the next T200 run the model's own wording contained no Trojan, counter or module names.
\end{itemize}

\section*{2026-10-09: LLM shown gates the detector missed (\texttt{cdc0acd})}
\begin{itemize}
  \item Flagged gates could fill the whole LLM pool: on T300 all 306 slots were flagged gates, so the LLM only re-confirmed the detector.
  \item Flagged gates now take at most 25\% of the pool (minimum 20); the other views skip flagged gates. On T300 the model sees 230 unflagged gates, including one payload gate the detector missed.
\end{itemize}

\section*{2026-10-09: free-running timers; reproducible simulation (\texttt{8ce6325})}
\begin{itemize}
  \item T300's 50-cycle window counter (32 registers on the ordinary clock, next state depends only on itself) matched no rule: 122 Trojan gates missed, recall 0.937.
  \item New evidence \texttt{autonomous\_timer} (lead, 0.60) / \texttt{trojan\_timer} (0.80 when it feeds other flagged logic). Zero matches on \textasciitilde{}50 clean ISCAS'89 netlists. Full T300 run afterwards: recall 0.937 $\rightarrow$ 0.9974.
  \item Random simulation iterated a set of inputs, so never-toggling counts changed between identical runs. Inputs are now sorted.
\end{itemize}

\section*{Earlier history (2026-07-24 to 2026-10-09, 128 commits)}
Summarised by phase from the commit history.

\subsection*{Phase 1: framework and deployment (Jul 24 to Aug 11)}
\begin{itemize}
  \item Initial PyATPG framework with a FastAPI backend and React frontend.
  \item Hosted on Render (backend) and Vercel (frontend). Atalanta is compiled from source and statically linked at build time, after glibc mismatches between Render's build and runtime containers.
  \item Bench format import and export; Atalanta failures reported loudly instead of as 0\% coverage; real CPU time and memory on Linux.
  \item First LLM ranker, a Trojan insertion tool, and removal of legacy ATPG code.
\end{itemize}

\subsection*{Phase 2: sequential ATPG and honest comparisons (Aug 14 to Aug 30)}
\begin{itemize}
  \item Sequential ATPG engine with a context-reduced LLM fault prioritiser, and its results in the frontend.
  \item Compare mode made to actually compare deterministic vs LLM-assisted runs.
  \item Gate names anonymised before the targeting prompt (\texttt{c23c039}).
  \item Cyclic netlists stopped before reaching Atalanta; false "feedback loop" errors on dangling signals fixed.
\end{itemize}

\subsection*{Phase 3: reliable Trojan generation and evaluation (Sep 2 to Sep 28)}
\begin{itemize}
  \item Trojan placements sampled instead of repeating; a trigger accepted only if it can actually activate; two variants never share a payload target.
  \item Redundant faults told apart from aborted ones.
  \item Multi-provider LLM support, unattended batch automation, a research pipeline, and zero-install runs for collaborators against the hosted API.
  \item Targeted coverage made honest: effort escalation, whole-circuit pattern simulation, compact LLM context, per-request token budgets.
  \item Exact parsing of synthesized hierarchical Verilog; a full-scan model for circuits with flip-flops; MUX payloads; an LLM judge for trigger/payload candidates with baselines.
\end{itemize}

\subsection*{Phase 4: Trust-Hub designs and the Trojan detector (Oct 1 to Oct 9)}
\begin{itemize}
  \item Up to five rotating LLM keys; Groq \texttt{openai/gpt-oss-120b} as the default.
  \item RTL designs accepted via Yosys synthesis; size and timeout guards for the hosted backend.
  \item Server stays responsive during runs, with live progress, cancellation of disconnected runs, an admin page and a PostgreSQL run database.
  \item Logic-clocked registers modelled (Trust-Hub sequential triggers); trigger detector; scan enable held low under full scan; LLM batches that carry a shortlist forward.
  \item Payload splices, hand-edited gates, evidence families that do not double count, region completion aiming for every Trojan gate, and recall reported with and without the LLM.
\end{itemize}

\end{document}
